\documentclass[addpoints,12pt]{exam}
%\documentclass[answers,addpoints,12pt]{exam}
\unframedsolutions

\renewcommand{\solutiontitle}{\noindent\textbf{}\enspace}
\SolutionEmphasis{\color{blue}}

\usepackage{empheq}
\usepackage{multirow}
%\usepackage{amsmath}
\usepackage[export]{adjustbox}
\usepackage{amsmath}%
\usepackage{MnSymbol}%
\usepackage{wasysym}%

\usepackage{setspace}
\usepackage{lastpage}
\usepackage{fullpage}
\usepackage{sectsty}
\usepackage{here}
\usepackage{hyperref}
%\usepackage{fancyhdr}
\usepackage{graphicx}
\usepackage[]{geometry}
\usepackage{color}
\usepackage{paralist}
%\setlength\textwidth{6.5in}
\geometry{paperwidth=8.5in,paperheight=11in,nohead,top=2cm,bottom=4cm, left=2.5cm, right=4cm}
\usepackage{textcomp}
\usepackage{array}
\usepackage{wrapfig}

\usepackage{ulem}

\newenvironment{packed_itemize}{
\begin{itemize}
\vspace{-12pt}
  \setlength{\topsep}{0pt}
  \setlength{\partopsep}{0pt}
  \setlength{\itemsep}{1pt}
  \setlength{\parskip}{0pt}
  \setlength{\parsep}{0pt}
}{\end{itemize}}


\newif\ifisAnswers
%\newif\ifisWhiteSpace

%\isAnswerstrue
\isAnswersfalse


% \newcommand{\department}[1]{\def\@department{#1}}
% 1 = ME
% 2 = EE
% 3 = Materials Science
% \usepackage{mciteplus}
% \usepackage{chapterbib}

\usepackage[numbers]{natbib}
\renewcommand{\bibnumfmt}[1]{#1.}
\renewcommand{\bibsection}{}

% \usepackage{multibib}[resetlabels]
\usepackage[resetlabels]{multibib}
\usepackage{multibib}
\newcites{ltex,soft,pres}%
    {Journal Papers,%
    Software,%
      Presentations}

% \newcites{papers}
% {Journal Papers}
% \newcites{papers,software,presentation}
% {Journal Papers,%
% Software,%
% Presentations}

\usepackage[compact]{titlesec}
\titlespacing{\subsection}{0pt}{*0}{*0}

\setlength{\headheight}{27pt}

\lhead{\normalsize ENGR0020 Probability \& Statistics for Engineers \\ 
University of Pittsburgh\\ \hrule}

%\lhead{ENGR 1066 Introduction to Solar Cells and Nanotechnology\\
%University of Pittsburgh\\ \hrule}


%\chead{\textbf{Assignment 1}}
%\rhead{\textbf{Due November 6, 2013}}
%\runningheadrule
\footer{}{Page \thepage\ of \numpages}{}

\checkedchar{\textcolor{blue}{X}}
\CorrectChoiceEmphasis{\color{blue}}

\newcommand{\red}[1]{\textcolor{red}{[#1]}} %for displaying red texts
%\newcommand{\red}[1]{} %for displaying red texts
\newcommand{\blue}[1]{\textcolor{blue}{#1}} %for displaying red texts
%\newcommand{\blue}[1]{} %for displaying red texts


% change the blue command, so leaves vertical white space
% 

% \newcommand{\HRule}{\rule[-.2mm]{13cm}{0.2mm}}
\newcommand{\sectionline}{%
  \nointerlineskip \vspace{\baselineskip}%
  \hspace{\fill}\rule{0.5\linewidth}{.7pt}\hspace{\fill}%
  \par\nointerlineskip \vspace{\baselineskip}
}
% \chead{Middle top}


\def\endncolumn{\end{tabular*}}

% \renewcommand{\labelitemi}{}


\checkboxchar{\underline{\hspace{1cm}}}

\begin{document}
% \department{1}
\headsep = 0pt

\sectionfont{\large}

%\singlespace 
\onehalfspace 

\vspace{0pt}

%{\noindent}{\textbf{IE 1051 Engineering Product Design}}\\
%{\noindent}{\textbf{Exam 2013, Written Portion}}
%\vspace{.1in}
%\hrule
%\vspace{0.1in}

%\pagestyle{head,foot}
\pagestyle{headandfoot}

%Define points
\pointsinrightmargin
\setlength{\rightpointsmargin}{0cm}
\marginpointname{ \points}

%\vspace{12pt}
{\noindent}{\textbf{Assignment \#10}}\\
%\vspace{12pt}
%\\
%{\noindent}Out: Wednesday, August 28, 2013\\
%{\noindent} Due: Tuesday, April 5, 2016

\begin{questions}
\singlespacing

%\question 
%Summary statistics for heights of fathers and sons are as follows:\\
%average height of fathers $\approx$ 68 inches, SD $\approx$ 2.5 inches\\
%average height of sons $\approx$ 69 inches, SD $\approx$ 2.7 inches, r $\approx$ 0.5\\
%(a) What is the relationship between 
%(a) Find the regression equation for predicting the height of a son from the
%height of his father.\\
%(b) Find the regression equation for predicting the height of a father from the
%height of his son.
%\begin{solution}
%$S_xx = \sum_i^{n} (x_i - \bar{x} )^2$
%$SD = \sqrt{S_{xx}}$
%$SD = \sqrt{S_{xx}}$
%r = $b_1 \times \sqrt{\frac{S_xx}{S_yy}}$
%$b_1 = $
%
%(a) the slope = $0.5 \times 2.7/2.7$ = 0.5\\
%For $-68 \times 0.5 = -34.$
%This would be $69 - 34 = 35$
%The equation is 
%\begin{equation*}
%\textrm{Height of son} = 0.5 \times \textrm{Height of father} + 35
%\end{equation*} \\
%(b) the slope is the same\\
%$For -69 \times -.5 = -34.5 $
%This would be $68 - 34.5 = 33.5$\\
%The equation is
%\begin{equation*}
%\textrm{Height of father} = 0.5 \times \textrm{Height of son} + 33.5
%\end{equation*} \\
%\end{solution}

\question Here is a hypothetical dataset called \texttt{height\_weight.csv}.
(a) Plot the
scatter diagram and draw the regression line for y on x.
\begin{solution}
\includegraphics[scale=.5]{height_weight.png}\\
%\includegraphics[scale=.25]{Ch10C4b}\\
%\includegraphics[scale=.25]{Ch10C4c}\\
%\includegraphics[scale=.25]{Ch10C4d}\\
\end{solution}
(b) What is the R$^2$?  What does this mean?  
\begin{solution}
.988.  It is the fraction of variability in $y$ explained by the regression model.  
\end{solution}




\question 
A small data set is shown below.
\begin{table}[H]
\begin{center}
\begin{tabular}{|c|c|}
\hline \hline
x & y \\ \hline \hline
1 & 2\\
2 & 3\\ 
3 & 1\\ 
4 & 5\\ 
5 & 6\\
\hline
\end{tabular}
\end{center}
\label{default}
\end{table}%
The fitted regression line is 
$y = x + .4$.  
\begin{parts}
\part What happens to the regression line if you  add 3 to each value of y instead of interchanging the
columns.
\begin{solution}
$y = x + 3.4$.  
The y-intercept is increased by 3, but the slope stays the same.  
\end{solution}
\part What happens to the regression line if you double each value in $x$.  
\begin{solution}
The slope is halved.  
$y = .5 x + .4$.  
\end{solution}
\end{parts}

\question The data file \texttt{couplesheights.txt} has been provided with the 
height of the husband in the first column in cm and the height of the wife in the second column in cm.  
\begin{parts}
\part Plot a scatter diagram.
\part Plot the fitted regression line of the height of the wife as a function of the height of the husband.  
\part Now plot the fitted regression line of the height of the husband as a function of the height of the wife.
Clearly label each line. 
\part Are the two regression lines the same or different?  Why?
\begin{solution} 
\includegraphics[scale=.5]{couplesHeights}
The two regression lines are different, because one regression line picks off the centers of the vertical strips
and the other picks off the center of the horizontal strips.  
\end{solution}
\end{parts}


\question The regression method predicted that a student at the 90th
percentile on the SAT would only be at the 69th percentile on first-year GPA. \\
True or false, and explain: a student at the 69th percentile on first-year GPA should be
at the 90th percentile on the SAT.  (Hint: Look at the previous question)
\begin{parts}
\begin{solution}
False.  The regression lines are different as shown from the previous questions.  
\end{solution}
\part Someone who scored at the 30th percentile on the SAT on average should score 
\underline{\hspace{1cm}} the 30th percentile on his/her first-year GPA.\\
Options: lower than, about equal to, higher than \\
Why?
\begin{solution}
higher than.  Regression to the mean.  
\end{solution}
\end{parts}



%\question 
%A professor in the School of Business in a university
%polled a dozen colleagues about the number of
%professional meetings they attended in the past five
%years ($x$) and the number of papers they submitted
%to refereed journals ($y$) during the same period. The
%summary data are given as follows:
%$n = 12$, $\bar{x} = 4$, $\bar{y} = 12$, 
%$\sum_{i = 1}^{n} x_i^2 = 232$, \\
%$\sum_{i = 1}^{n} x_i y_i = 318$, 
%Fit a simple linear regression model between x and y by
%finding out the estimates of intercept and slope. Comment
%on whether attending more professional meetings
%would result in publishing more papers.
%\begin{solution}
%\begin{align*}
%b_1 &= \frac{n \sum_{i = 1}^{n} x_i y_i  - n \bar{x} n \bar{y}}{n \sum_{i = 1}^{n} x_i^2 - ( n \bar{x} )^2}  \\
%&= \frac{12 \times 318 - 12 \times 4 \times 12 \times 12}{12 \times 232 - (12 \times 4)^2} \\
%&= -6.45
%\end{align*}
%\begin{align*}
%b_0 &= \bar{y} - b_1 \bar{x} \\
%&= 12 + 6.45 \times 4 = 37.8
%\end{align*}
%Attending more professional meetings results in publishing less papers.  
%\end{solution}

\question 
A study was made on the amount of converted
sugar in a certain process at various temperatures. 
The data is given in the file, \texttt{temperature\_sugar.txt}
with the temperature in the first column and the converted sugar in the second column.  \\
%The
%data were coded and recorded as follows:
%\begin{table}[H]
%%\caption{default}
%\begin{center}
%\begin{tabular}{|c|c|}
%\hline \hline
%Temperature, x & Converted Sugar, y \\ \hline \hline
%1.0 & 8.1 \\
%1.1 & 7.8 \\
%1.2 & 8.5 \\
%1.3 & 9.8 \\
%1.4 & 9.5 \\
%1.5 & 8.9 \\
%1.6 & 8.6 \\
%1.7 & 10.2 \\
%1.8 & 9.3 \\
%1.9 & 9.2 \\
%2.0 & 10.5 \\
%\hline \hline
%\end{tabular}
%\end{center}
%\label{default}
%\end{table}%
(a) Estimate the linear regression line.\\
(b) Estimate the mean amount of converted sugar produced
when the coded temperature is 1.75.\\
(c) Plot the residuals versus temperature. Comment.
\begin{solution}
(a) $y = 1.8 x + 6.4$\\
(b) $9.6$\\
(c)
\includegraphics[scale=.5]{residuals}
Residuals appear to be random as desired.  

\end{solution}

\question For the previous question,
(a) evaluate $s^2$;
(b) construct a 95\% confidence interval for $\beta_0$;
(c) construct a 95\% confidence interval for $\beta_1$.
\begin{solution}
(a) $s^2 = .4$
(b) $6.4136 \pm 2.0917$
(c) $1.809 \pm 1.36$
\end{solution}

\question 
A mathematics placement test is given to all entering
freshmen at a small college. A student who receives
a grade below 35 is denied admission to the regular
mathematics course and placed in a remedial class.
The placement test scores and the final grades for 20
students who took the regular course were recorded.
The data is given in file, \texttt{placement\_grade.txt} 
with the placement test scores in the first column and the course grade scores in the second column.  \\
(a) Plot a scatter diagram.\\
(b) Find the equation of the regression line to predict
course grades from placement test scores.\\
(c) Graph the line on the scatter diagram.\\
(d) If 60 is the minimum passing grade, below which
placement test score should students in the future
be denied admission to this course?
%\begin{table}[H]
%%\caption{default}
%\begin{center}
%\begin{tabular}{|c|c|} \hline
%Placement Test & Course Grade \\ \hline
%50 & 53 \\
%35 & 41 \\
%35 &61 \\
%40 &56\\
%55 &68\\
%65 &36\\
%35 &11\\
%60 &70\\
%90 &79\\
%35 &59\\
%90 &54\\
%80 &91\\
%60 &48\\
%60 &71\\
%60 &71\\
%40 &47\\
%55 &53\\
%50 &68\\
%65 &57\\
%50 &79\\ \hline
%\end{tabular}
%\end{center}
%\label{default}
%\end{table}
\begin{solution}
(a) \includegraphics[scale=.3]{p11p8}\\
(b) y = .471 x + 32.5\\
(d) 58.4\\
\end{solution}

\question 
Test the hypothesis that $\beta_0 = 10$ in the previous problem against the alternative that $\beta_0 < 10$.
Use a 0.05 level of significance.
\begin{solution}
%$beta0_0 = 10$
%$t_stat = (beta0_hat - beta0_0) / se_beta0$
%p_value = stats.t.cdf(t_stat, df_resid)

$H_0: \beta_0 = 10$\\
$H_0: \beta_0 < 10$\\
$t = 1.78$\\
$P = .95$\\
Fail to reject $H_0$
\end{solution}

\question Six different machines are being considered for
use in manufacturing rubber seals. The machines are
being compared with respect to tensile strength of the
product. A random sample of four seals from each machine
is used to determine whether the mean tensile
strength varies from machine to machine. The following
are the tensile-strength measurements in kilograms
per square centimeter $\times 10^{-1}$:
\begin{table}[H]
%\caption{default}
\begin{center}
\begin{tabular}{|c|c|c|c|c|c|}
\hline
1&2&3&4&5&6 \\ \hline
17.5&16.4&20.3&14.6&17.5&18.3 \\
16.9&19.2&15.7&16.7&19.2&16.2 \\
15.8&17.7&17.8&20.8&16.5&17.5 \\
18.6&15.4&18.9&18.9&20.5&20.1 \\ \hline
\end{tabular}
\end{center}
\label{default}
\end{table}%
Perform the analysis of variance at the 0.05 level of significance
and indicate whether or not the mean tensile
strengths differ significantly for the six machines.
\begin{solution}
$H_0: \mu_1 = \mu_2 = \ldots \mu_6$\\
$H_1:$ At least two of the means are not equal\\
Use the f-test.\\
$P = .90$
Do not reject $H_0$\\
\end{solution}


\question The data in the following table represent the
number of hours of relief provided by five different
brands of headache tablets administered to 25 subjects
experiencing fevers of 38$^{\circ}$
C or more. Perform the analysis
of variance and test the hypothesis at the 0.05 level
of significance that the mean number of hours of relief
provided by the tablets is the same for all five brands.
Discuss the results.
\begin{table}[H]
%\caption{default}
\begin{center}
\begin{tabular}{|c|c|c|c|c|c|}
\hline
A&B&C&D&E\\ \hline
5.2&9.1&3.2&2.4&7.1\\
4.7&7.1&5.8&3.4&6.6\\
8.1&8.2&2.2&4.1&9.3\\
6.2&6&3.1&1&4.2\\
3&9.1&7.2&4&7.6\\ \hline
\end{tabular}
\end{center}
\label{default}
\end{table}%
\begin{solution}
$H_0: \mu_1 = \mu_2 = \ldots \mu_5$\\
$H_1:$ At least two of the means are not equal\\
Use the f-test.\\
$P = .0015$\\
Reject $H_0$
\end{solution}



%\question Six different machines are being considered for
%use in manufacturing rubber seals. The machines are
%being compared with respect to tensile strength of the
%product. A random sample of four seals from each machine
%is used to determine whether the mean tensile
%strength varies from machine to machine. The following
%are the tensile-strength measurements in kilograms
%per square centimeter $\times 10^{-1}$:
%\begin{table}[H]
%%\caption{default}
%\begin{center}
%\begin{tabular}{|c|c|c|c|c|c|}
%\hline
%1&2&3&4&5&6 \\ \hline
%17.5&16.4&20.3&14.6&17.5&18.3 \\
%16.9&19.2&15.7&16.7&19.2&16.2 \\
%15.8&17.7&17.8&20.8&16.5&17.5 \\
%18.6&15.4&18.9&18.9&20.5&20.1 \\ \hline
%\end{tabular}
%\end{center}
%\label{default}
%\end{table}%
%Perform the analysis of variance at the 0.05 level of significance
%and indicate whether or not the mean tensile
%strengths differ significantly for the six machines.
%\begin{solution}
%$H_0: \mu_1 = \mu_2 = \ldots \mu_6$\\
%$H_1:$ At least two of the means are not equal\\
%Use the f-test.\\
%$P = .90$
%Do not reject $H_0$\\
%\end{solution}
%
%
%\question 13.2  The data in the following table represent the
%number of hours of relief provided by five different
%brands of headache tablets administered to 25 subjects
%experiencing fevers of 38$^{\circ}$
%C or more. Perform the analysis
%of variance and test the hypothesis at the 0.05 level
%of significance that the mean number of hours of relief
%provided by the tablets is the same for all five brands.
%Discuss the results.
%\begin{table}[H]
%%\caption{default}
%\begin{center}
%\begin{tabular}{|c|c|c|c|c|c|}
%\hline
%A&B&C&D&E\\ \hline
%5.2&9.1&3.2&2.4&7.1\\
%4.7&7.1&5.8&3.4&6.6\\
%8.1&8.2&2.2&4.1&9.3\\
%6.2&6&3.1&1&4.2\\
%3&9.1&7.2&4&7.6\\ \hline
%\end{tabular}
%\end{center}
%\label{default}
%\end{table}%
%\begin{solution}
%$H_0: \mu_1 = \mu_2 = \ldots \mu_5$\\
%$H_1:$ At least two of the means are not equal\\
%Use the f-test.\\
%$P = .0015$\\
%Reject $H_0$
%\end{solution}
%
%
%
%
%\question 13.3  In an article ``Shelf-Space Strategy in Retailing,Ó
%published in Proceedings: Southern Marketing Association,
%the effect of shelf height on the supermarket sales
%of canned dog food is investigated. An experiment was
%conducted at a small supermarket for a period of 8 days
%on the sales of a single brand of dog food, referred to
%as Arf dog food, involving three levels of shelf height:
%knee level, waist level, and eye level. During each day,
%the shelf height of the canned dog food was randomly
%changed on three different occasions. The remaining
%sections of the gondola that housed the given brand
%were filled with a mixture of dog food brands that were
%both familiar and unfamiliar to customers in this particular
%geographic area. Sales, in hundreds of dollars,
%of Arf dog food per day for the three shelf heights are
%given. Based on the data, is there a significant difference
%in the average daily sales of this dog food based
%on shelf height? Use a 0.01 level of significance.
%\begin{table}[H]
%%\caption{default}
%\begin{center}
%\begin{tabular}{|c|c|c|c|c|c|}
%\hline
%Knee Level & Waist Level & Eye Level \\ \hline
%77&88&85\\
%82&94&85\\
%86&93&87\\
%78&90&81\\
%81&91&80\\
%86&94&79\\
%77&90&87\\
%81&87&93\\ \hline
%\end{tabular}
%\end{center}
%\label{default}
%\end{table}%
%\begin{solution}
%$H_0: \mu_1 = \mu_2 = \ldots \mu_3$\\
%$H_1:$ At least two of the means are not equal\\
%Use the f-test.\\
%$P = .00011$\\
%Reject $H_0$\\
%\end{solution}
%
%
%\question 13.5   The mitochondrial enzyme NADPH:NAD
%transhydrogenase of the common rat tapeworm (Hymenolepiasis
%diminuta) catalyzes hydrogen in the
%transfer from NADPH to NAD, producing NADH.
%This enzyme is known to serve a vital role in the
%tapewormÕs anaerobic metabolism, and it has recently
%been hypothesized that it may serve as a proton exchange
%pump, transferring protons across the mitochondrial
%membrane. A study on Effect of Various
%Substrate Concentrations on the Conformational Variation
%of the NADPH:NAD Transhydrogenase of Hymenolepiasis
%diminuta, conducted at Bowling Green
%State University, was designed to assess the ability of
%this enzyme to undergo conformation or shape changes.
%Changes in the specific activity of the enzyme caused
%by variations in the concentration of NADP could be
%interpreted as supporting the theory of conformational
%change. The enzyme in question is located in the inner
%membrane of the tapewormÕs mitochondria. Tapeworms
%were homogenized, and through a series of centrifugations,
%the enzyme was isolated. Various concentrations
%of NADP were then added to the isolated
%enzyme solution, and the mixture was then incubated
%in a water bath at 56$^{\circ}$C for 3 minutes. The enzyme
%was then analyzed on a dual-beam spectrophotometer,
%and the results shown were calculated, with the specific
%activity of the enzyme given in nanomoles per minute
%per milligram of protein. Test the hypothesis at the
%0.01 level that the average specific activity is the same
%for the four concentrations.
%
%\begin{table}[H]
%%\caption{default}
%\begin{center}
%\begin{tabular}{|c|c|c||c|c|}
%\hline
%0 & 80 & 160 & \multicolumn{2}{|c|}{360} \\ \hline
%11.01&11.38&11.02&6.04&10.31 \\
%12.09&10.67&10.67&8.65&8.3 \\
%10.55&12.33&11.5&7.76&9.48 \\
%11.26&10.08&10.31&10.13&8.89 \\
%&&&9.36& \\ \hline
%\end{tabular}
%\end{center}
%\label{default}
%\end{table}%
%\begin{solution}
%$H_0: \mu_1 = \mu_2 = \ldots \mu_4$\\
%$H_1:$ At least two of the means are not equal\\
%Use the f-test.\\
%$P = 0.022$\\
%Reject $H_0$\\
%\end{solution}
%
%%
%%\begin{solution}
%%
%%\end{solution}
%%
%%\question 13.9.  Use BartlettÕs test at the 0.01 level of significance
%%to test for homogeneity of variances in Exercise
%%13.5 on page 519.
%%\begin{solution}
%%$H_0: \sigma_1^2 = \sigma_2^2 = \sigma_3^2 = \sigma_4^2$\\
%%$H_1:$ The variances are not equal\\
%%\end{solution}
%%
%%\question 13.10
%%Use CochranÕs test at the 0.01 level of significance
%%to test for homogeneity of variances in Exercise
%%13.4 on page 519.
%%
%%\question 13.12
%%Consider the data of Review Exercise 13.45 on
%%page 555. Make significance tests on the following contrasts:
%%(a) B versus A, C, and D;
%%(b) C versus A and D;
%%(c) A versus D.
%%\begin{solution}
%%\begin{table}[H]
%%%\caption{default}
%%\begin{center}
%%\begin{tabular}{|c|c|c|c|c|}
%%\hline
%%Source of Variation & Sum of Squares & Degrees of Freedom & Mean Square & Computed f \\ \hline
%%B vs.~A, C, D & 30.67 & 1 & 30.67 & 14.28 \\
%%C vs.~A, D & 49.92 & 1 & 49.92 & 23.23 \\
%%A vs.~D & 5.33 & 1 & 5.33 & 2.48 \\
%%Error & 34.38 & 16 & 2.15 & \\ \hline
%%\end{tabular}
%%\end{center}
%%\label{default}
%%\end{table}%
%%(a) P-value = 0.0016; different\\
%%(b) P-value = 0.0002; different\\
%%(c) P-value = 0.1349; same\\
%%\end{solution}
%
%%\question 13.13
%%The purpose of the study The Incorporation
%%of a Chelating Agent into a Flame Retardant Finish
%%of a Cotton Flannelette and the Evaluation of Selected
%%Fabric Properties conducted at Virginia Tech was to
%%evaluate the use of a chelating agent as part of the
%%flame-retardant finish of cotton flannelette by determining
%%its effects upon flammability after the fabric is
%%laundered under specific conditions. Two baths were
%%prepared, one with carboxymethyl cellulose and one
%%without. Twelve pieces of fabric were laundered 5 times
%%in bath I, and 12 other pieces of fabric were laundered
%%10 times in bath I. This was repeated using 24 additional
%%pieces of cloth in bath II. After the washings the
%%lengths of fabric that burned and the burn times were
%%measured. For convenience, let us define the following
%%treatments:
%%
%%
%%Treatment 1: 5 launderings in bath I,\\
%%Treatment 2: 5 launderings in bath II,\\
%%Treatment 3: 10 launderings in bath I,\\
%%Treatment 4: 10 launderings in bath II.\\
%%
%%Burn times, in seconds, were recorded as follows:
%%
%%\begin{table}[H]
%%%\caption{default}
%%\begin{center}
%%\begin{tabular}{|c|c|c|c|}
%%\hline
%%\textbf{A} & \textbf{B} & \textbf{C} & \textbf{D} \\
%%13.7&6.2&27.2&18.2\\
%%23&5.4&16.8&8.8\\
%%15.7&5&12.9&14.5\\
%%25.5&4.4&14.9&14.7\\
%%15.8&5&17.1&17.1\\
%%14.8&3.3&13&13.9\\
%%14&16&10.8&10.6\\
%%29.4&2.5&13.5&5.8\\
%%9.7&1.6&25.5&7.3\\
%%14&3.9&14.2&17.7\\
%%12.3&2.5&27.4&18.3\\
%%12.3&7.1&11.5&9.9\\ \hline
%%\end{tabular}
%%\end{center}
%%\label{default}
%%\end{table}%
%%
%%
%% 
%%(a) Perform an analysis of variance, using a 0.01 level of
%%significance, and determine whether there are any
%%significant differences among the treatment means.
%%\begin{solution}
%%$H_0: \mu_1 = \mu_2 = \ldots \mu_4$\\
%%$H_1:$ At least two of the means are not equal\\
%%$P = 2.23$\\
%%Reject $H_0$ and accept $H_1$\\
%%\end{solution}
%%(b) Use single-degree-of-freedom contrasts with $\alpha =
%%0.01$ to compare the mean burn time of treatment
%%1 versus treatment 2 and also treatment 3 versus
%%treatment 4.
%%\begin{solution}
%%For treatment 1 versus 2, $P = 0.0000024$
%%They are different\\
%%For treatment 1 versus 2, $P = 0.0648$
%%Accept they are the same\\
%%\end{solution}
%%
%%
%%
%%
%%
%%
%%\question 14.1
%%An experiment was conducted to study the effects
%%of temperature and type of oven on the life of
%%a particular component. Four types of ovens and
%%3 temperature levels were used in the experiment.
%%Twenty-four pieces were assigned randomly, two to
%%each combination of treatments, and the following results
%%recorded.
%%
%%\begin{table}[H]
%%%\caption{default}
%%\begin{center}
%%\begin{tabular}{|c|c|c|c|c|}
%%\hline \hline
%%\textbf{Temperature ($^{\circ}$ F)} & \textbf{O$_1$} & \textbf{O$_2$} & \textbf{O$_3$} & \textbf{O$_4$} \\ \hline \hline
%%\textbf{500} & 227&214&225&260 \\
%%& 221&259&236&229 \\ \hline
%%\textbf{550} & 187&181&232&246 \\
%%& 208&179&198&273 \\ \hline
%%\textbf{600} & 174&198&178&206 \\
%%& 202&194&213&219 \\ \hline \hline
%%\end{tabular}
%%\end{center}
%%\label{default}
%%\end{table}%
%%
%%Using a 0.05 level of significance, test the hypothesis
%%that
%%(a) different temperatures have no effect on the life of
%%the component;
%%(b) different ovens have no effect on the life of the component;
%%(c) the type of oven and temperature do not interact.
%%\begin{solution}
%%(a) $P = .0059$; Reject $H_0$\\
%%(b) $P = .0159$; Reject $H_0$\\
%%(c) $P = .2215$; Do not reject $H_0$
%%\end{solution}
%%
%%\question 14.2
%%To ascertain the stability of vitamin C in reconstituted
%%frozen orange juice concentrate stored in a
%%refrigerator for a period of up to one week, the study
%%Vitamin C Retention in Reconstituted Frozen Orange
%%Juice was conducted by the Department of Human Nutrition
%%and Foods at Virginia Tech. Three types of
%%frozen orange juice concentrate were tested using 3 different
%%time periods. The time periods refer to the number
%%of days from when the orange juice was blended
%%until it was tested. The results, in milligrams of ascorbic
%%acid per liter, were recorded. Use a 0.05 level of
%%significance to test the hypothesis that
%%(a) there is no difference in ascorbic acid contents
%%among the different brands of orange juice concentrate;
%%(b) there is no difference in ascorbic acid contents for
%%the different time periods;
%%(c) the brands of orange juice concentrate and the
%%number of days from the time the juice was blended
%%until it was tested do not interact.
%%
%%\begin{table}[H]
%%%\caption{default}
%%\begin{center}
%%\begin{tabular}{|c|c|c|c|c|c|c|}
%%\hline 
%%& \multicolumn{6}{|c|}{\textbf{Time (days)}} \\ \hline
%%\textbf{Brand} & \multicolumn{2}{|c|}{\textbf{0}} & \multicolumn{2}{|c|}{\textbf{3}} & \multicolumn{2}{|c|}{\textbf{7}} \\ \hline
%%\textbf{Richfood} & 52.6&54.2&49.4&49.2&42.7&48.8 \\
%%& 49.8&46.5&42.8&53.2&40.4&47.6 \\ \hline
%%\textbf{Sealed-Sweet} & 56.0&48.0&48.8&44.0&49.2&44 \\
%%& 49.6&48.4&44.0&42.4&42.0&43.2 \\ \hline
%%\textbf{Minute Maid} & 52.5&52.0&48.0&47.0&48.5&43.3 \\
%%&51.8&53.6&48.2&49.6&45.2&47.6 \\ \hline
%%\end{tabular}
%%\end{center}
%%\label{default}
%%\end{table}%
%%
%%\begin{solution}
%%(a) $P = .0059$; Reject $H_0$\\
%%(b) $P = .0159$; Reject $H_0$\\
%%(c) $P = .2215$; Do not reject $H_0$
%%\end{solution}
%
%\question 
%A study is made of incomes among full-time workers age 25-54 in a certain town.
%A simple random sample is taken, of 250 people with high school degrees: the
%sample average income is \$30,000 and the SD is \$25,000. Another simple random
%sample is taken, of 250 people with college degrees: the sample average income is
%\$50,000 and the SD is \$40,000.
%Question: Is the difference in averages real, or due to chance?
%To answer this question, which do you use?
%(i) the one-sample $z$-test.
%(ii) the two-sample $z$-test.
%(iii) the one-sample $t$-test.
%(iv) the two-sample $t$-test.
%(iv) the $\chi^2$-test
%%, with a null hypothesis that tells you the c
%%( section I ) .
%(v) the $F$-test.\\
%Now answer the question.
%\begin{solution}
%The two sample $z$-test.  
%The difference is $20000$.  The SE for the difference is $2983.3$.  
%The difference is definitely real.  
%It is over 6 SEs.  
%\end{solution}
%
%
%\question 
%Demographers think that about 55\% of newborns are male. In a certain hospital,
%568 out of 1,000 consecutive births are male.
%Question: Are the data consistent with the theory?
%To answer this question, which do you use?
%(i) the one-sample $z$-test.
%(ii) the two-sample $z$-test.
%(iii) the one-sample $t$-test.
%(iv) the two-sample $t$-test.
%(iv) the $\chi^2$-test
%%(iv) the $\chi^2$-test for independence (section 4).
%(v) the $F$-test.\\
%%, with a null hypothesis that tells you the c
%%( section I ) .
%
%
%Now answer the question.
%\begin{solution}
%(i) the one-sample $z$-test.\\
%Null-hypothesis; 55\% of newborns are male\\
%Alternative-hypothesis; 55\% of newborns are not male.\\
%$z = (568 - 550)/15.7321 = 1.14$
%Should be two-sided test.  
%$P(Z > 1.1442 or Z < 1.1442) = .25$
%Accept the null-hypothesis.  
%\end{solution}
%
%\question A gambler is accused of using a loaded die, but he pleads innocent.
%You are called in to interpret the data and determine if the die is fair.  
%Suppose the observed frequencies in table 2 had come out as shown in table 7A
%below. Compute the value of $\chi^2$, the degrees of freedom, and P. What can be
%inferred?
%\begin{solution}
%The degrees of freedom is 5. 
%$\chi^2 = 13.2$.\\
%P = .0216
%\end{solution}
%
%
%
%\begin{table}[H]
%%\caption{default}
%\begin{center}
%\begin{tabular}{ || c|c||c|c || c|c || c|c ||}
%\hline
%\multicolumn{2}{||c ||}{Table 7 A} & \multicolumn{2}{ c ||}{Table 7 B} &  \multicolumn{2}{ c ||}{Table 7 C} & \multicolumn{2}{ c ||}{Table 7 D}  \\ \hline
%& Observed & & Observed & & Observed & & Observed \\
%Value & frequency &  Value &  frequency &  Value &  frequency &  Value &  frequency \\
%1 & 5 & 1 & 9 & 1 & 90 & 1 & 10,287 \\
%2 & 7 & 2 & 11 & 2 & 110 & 2 & 10,056 \\
%3 & 17 & 3 & 10 & 3 & 100 & 3 & 9,708 \\
%4 & 16 & 4 & 8 & 4 & 80 & 4 & 10,080 \\
%5 & 8 & 5 & 12 & 5 & 120 & 5 & 9,935 \\
%6 & 7 & 6 & 10 & 6 & 100 & 6 & 9,934 \\ \hline
%\end{tabular}
%\end{center}
%\label{default}
%\end{table}%
%
%
%\question 
%Suppose the observed frequencies in table l had come out as shown in table 7B.
%Make a $\chi^2$-test of the null hypothesis that the die is fair.
%\begin{solution}
%The degrees of freedom is 5. 
%$\chi^2 = 1$\\
%P = .96\\
%The die is fair.
%\end{solution}
%
%
%\question 
%Suppose that table I had 600 entries instead of 60, with observed frequencies as
%shown in table 7C. Make a $\chi^2$-test of the null hypothesis that the die is fair.
%\begin{solution}
%The degrees of freedom is 5. 
%$\chi^2 = 10$\\
%$P = .0752$
%The die still appears fair.
%\end{solution}
%
%\question 
%Suppose that table 1 had 60,000 entries, with the observed frequencies as shown
%in table 7D.
%(a) Compute the percentage of times each value showed up.
%(b) Does the die look fair?
%(c) Make a $\chi^2$-test of the null hypothesis that the die is fair.
%\begin{solution}
%(a) .174, .168, .162, .168, 166, .166\\
%(b) Yes\\
%(c) $chi^2 = 18.6$;\\ 
%$P = .0023$; the die looks loaded
%\end{solution}
%


\end{questions}

\end{document} 